\section{本章小结}
\label{sec:saddle-summary}

本章从一个具体的疑问出发：梯度很小的点算不算找到了解。梯度小只说明一阶变化弱，山口和谷底都满足这一点。

驻点按二阶信息分成严格鞍点和二阶驻点。负曲率是确凿的证据，说明这里还能下降；Hessian正定也是确凿的证据，说明这里是局部极小点。半正定而奇异的情形，二阶信息无法判断，$x_1^2+x_2^4$与$x_1^2-x_2^4$就是例子。即使确实是局部极小点，也不说明是全局最优解。

负曲率可以只靠Hessian向量积发现。幂迭代和Lanczos方法从随机起点出发，找到方向时的结论是确定的，未找到时的结论带有失败概率，检测误差和差分误差都要计入。

离开鞍点有三种机制，它们的条件和结论各不相同。

\begin{BookInlineTable}
  \centering
  \begin{tabular}{@{}lll@{}}
    \toprule
    机制 & 需要的条件 & 结论\\
    \midrule
    负曲率步 & 方向、Hessian Lipschitz & 每步确定下降$\frac{2\gamma^3}{3M^2}$\\
    扰动梯度法 & 梯度、$L$、$M$、随机扰动 & 高概率到达近似二阶驻点\\
    不定模型 & 子问题使用负曲率 & 二阶必要条件的下极限\\
    \bottomrule
  \end{tabular}
  \caption{三种离开鞍点的机制。修正Hessian只为得到下降方向，它的保证停留在一阶。}
  \label{tab:saddle-summary}
\end{BookInlineTable}

检测精度和离开机制共同决定了能够声称的保证。停止时用两个容差描述近似二阶驻点，$\varepsilon_H=\sqrt{M\varepsilon_g}$是自然的搭配，它限制了小球内还能下降多少，但不说明这个点是局部极小点。计算量要把梯度、Hessian向量积、分解和扰动窗口分开统计，并报告确定性保证或相应的失败概率。

到这里，光滑优化中的方向、步长、曲率和成本的框架已经完整。这些方法都假设梯度存在。下一部分研究目标在某些点没有普通梯度时，怎样改用次微分和邻近结构。
